
\lesson{02/11}\footnote{
most of this lesson was a review of the previous online lesson
}
\textbf{Note} The contraction of two fields is a \textit{complex number}, for the
commutation relations of $a,a^\dagger$ (show it!), so
$$\boxed{T\left\{\phi(x)\phi(y)\right\} = N\left\{\phi(x)\phi(y)+
\overline{\phi(x)\phi(y)}\right\}}$$

With this result, the free propagator's expression greatly simplifies to
$$\bra{0}T\left\{\phi(x)\phi(y)\right\}\ket{0}=
\underbrace{\bra{0}N\left\{\phi(x)\phi(y)\right\}\ket{0}}_{:=0}+
\bra{0}\overline{\phi(x)\phi(y)}\ket{0} = \overline{\phi(x)\phi(y)}=
i\Delta_F(x-y)$$

The generalization of this result is called \textbf{Wick's theorem}:\\
\textbf{Wick's theorem}: \textit{The time ordered product of fields is
equal to the sum of the normal ordering of the product with all possible
contractions between the fields:}
$$T\left\{\prod_{i}\phi(x_i)\right\} = N\left\{\prod\phi_i(x_i)\right\}+
\sum_{ij}\overline{\phi(x_i)\phi(x_j)}$$


\textbf{Note} The contractions, as seen before, are all the possible Feynman
free propagators of the fields.
